\subsection{A uniform bound for abelian prime-to-$p$ monodromy}

\begin{proposition}[F5: uniform twisted cohomology in a geometric subclass; Proved]
Let $A$ be a complete normal two-dimensional local domain of
characteristic $p>0$ with algebraically closed residue field $k$,
and put $U=\Spec A\setminus\{\mathfrak m\}$. There is a constant
$C_A$ such that, for every finite locally constant
$\mathbf F_p$-sheaf $M$ on $U$ whose monodromy group is abelian
of order prime to $p$,
\[
 \dim_{\mathbf F_p}H^1_{\mathrm{et}}(U,M)
 \leq C_A\operatorname{rank}_{\mathbf F_p}M.
\]
Consequently the group in F3 cannot be a quotient of such a local
fundamental group when $p=2$.
\end{proposition}

\begin{proof}
We first produce one finite cover, independent of $M$, on whose
fixed resolution all these coefficient systems extend without
ramification. We then make Suzuki's equal-characteristic
nearby-cycle argument uniform for the resulting torsion line bundles.
The argument uses the explicit filtration and Frobenius-cokernel
calculation of \cite[Propositions 2.1 and 3.1--3.4, pp.~4--7]{Suzuki2023};
the uniformity and the preliminary fixed-cover reduction are proved below.

\paragraph{A finite bound for abelian tame inertia.}
Choose a coefficient field $k\hookrightarrow A$ by the equal-characteristic
Cohen structure theorem. Since $k$ is perfect, a power-series presentation
makes $A$, and every scheme of finite type over it used below, $F$-finite.
Choose a resolution $f:X\to\Spec A$ with reduced exceptional divisor
$D=\bigcup_{i=1}^s D_i$ a simple normal crossing divisor, each $D_i$
smooth and $D$ nonempty; if $A$ is regular, first blow up its closed point.
Its intersection matrix $B=(D_i\cdot D_j)$ is negative definite,
so $b=|\det B|$ is a positive integer. Write
$Q=\pi_1(U)^{\mathrm{ab},(p')}$ and let $I_i\subset Q$ be the image
of tame inertia at $D_i$. These are procyclic groups.

For every character $\chi:Q\to\mu_n$ with $(n,p)=1$, Kummer theory
represents the torsor by a line bundle $L$ on $U$ with a trivialization
of $L^{\otimes n}$. Extend $L$ to a line bundle $\widetilde L$ on the
regular scheme $X$: extend a representing divisor and use regularity.
The trivialization identifies
\[
 \widetilde L^{\otimes n}\simeq
 \mathcal O_X\Bigl(\sum_i a_iD_i\Bigr)
\]
for integers $a_i$. Its inertia character at $D_i$ is $a_i\bmod n$,
with the usual common choice of roots of unity. Restriction to $D_j$
gives
\[
 \sum_i a_i(D_i\cdot D_j)
       =n\deg(\widetilde L|_{D_j}).
\]
The adjugate matrix therefore gives $ba_i=0\bmod n$ for every $i$.
Characters separate points of the abelian pro-prime-to-$p$ group $Q$,
so $bI_i=0$. The closed subgroup $I=\sum_i I_i$ is finite.

The same subgroup contains the image of tame inertia for every
exceptional divisorial valuation over $\Spec A$. Indeed, at the center
of that valuation on $X$, trivialize $\widetilde L$. The Kummer equation
has the form
\[
 z^n=u\prod_{i\in S}t_i^{a_i},
\]
where $u$ is a unit and the $t_i$ are the local equations of the
components through the center. Its valuation is
$\sum_{i\in S}v(t_i)a_i$. The unit contributes only unramified
monodromy. In particular, if $\chi$ is trivial on $I$, then all $a_i$ are zero
modulo $n$, so this valuation inertia also has trivial character.
Characters of $Q/I$ separate points; this places its image in $I$
without making compatible choices of inertia generators.
There is no wild inertia in an abelian prime-to-$p$ quotient.

Choose an open subgroup $Q'\subset Q$ with $Q'\cap I=0$; a finite
quotient separating the finitely many nonzero elements of $I$ does
this. Let $U'\to U$ be the connected finite abelian prime-to-$p$
cover defined by $Q'$. Normalize $\Spec A$ in its function field to
obtain $\Spec A'$; excellence makes this normalization finite, and
completeness and the algebraically closed residue field make $A'$
a complete normal local domain with residue field $k$.
Choose an SNC resolution $X'\to\Spec A'$ with reduced exceptional
curve $D'$. For every $M$ in the proposition, its pullback to $U'$
has trivial inertia at every exceptional divisor of $X'$: its image
in $Q$ lies both in $I$ and in $Q'$. Here a divisorial valuation of
the finite field extension restricts to an exceptional divisorial
valuation of the base field, because residue transcendence degree one
is preserved under finite extension. Normalize $X'$ in its finite
monodromy cover. It is unramified in codimension one, hence is
\'etale everywhere by purity of the branch locus on the regular
scheme $X'$ \cite[Lemma 58.21.4, tag 0BMB]{StacksBranchPurity}.
Thus every such $M|_{U'}$ extends to a locally constant
sheaf on the same fixed $X'$.

Restriction
$H^1(U,M)\to H^1(U',M|_{U'})$ is injective, since restriction followed
by corestriction is multiplication by the degree, which is prime to
$p$. It suffices to obtain a uniform bound on $U'$. In the rest of
the proof rename $A',X',D',U'$ as $A,X,D,U$; these are now fixed.

\paragraph{The extended coefficient bundle.}
Let $M$ extend to $X$, set $r=\operatorname{rank}_{\mathbf F_p}M$, and
let $V=M\otimes_{\mathbf F_p}\mathcal O_X$. This is a vector bundle
with its natural Frobenius-semilinear map $F_V$, whose linearization
$F_X^*V\to V$ is an isomorphism. Because the monodromy is finite,
abelian, and of order prime to $p$, the representation over $k$
splits into characters. Hence
\[
 V\simeq\bigoplus_{a=1}^r L_a,
 \qquad L_a\text{ a finite-order line bundle on }X.
\]
In particular $\deg(L_a|_{D_i})=0$ for every $a,i$.
The summands need not be individually Frobenius-stable.

Choose a positive exceptional cycle $Z=\sum_i c_iD_i$ such that
$J=\mathcal O_X(-Z)$ has positive degree on every $D_i$, as in
\cite[p.~6]{Suzuki2023}. All choices below depend on this fixed
resolution and $J$, not on $M$.
With $j:U\hookrightarrow X$, $i:D\hookrightarrow X$, and
$\Psi=i^*j_*$ on \'etale sheaves, the Artin--Schreier sequence for
$V$ gives
\[
 0\longrightarrow V(*D)/V
 \xrightarrow{F_V-1} V(*D)/V
 \longrightarrow R^1\Psi M\longrightarrow0.
\]
Here we identify coherent quotients supported on $D$ with their
pullbacks by $i$. Exactness follows locally from the ordinary
Artin--Schreier sequence after an \'etale trivialization of $M$;
$j$ is affine and $F_V-1$ is surjective on the extended bundle in
the \'etale topology.

Let $U_mR^1\Psi M$ be the image of $J^{-pm}V/V$.
Valuations along the regular components of $D$, and the fact that
$F_X^*V\to V$ is an isomorphism, give the exact sequence
\[
 0\longrightarrow J^{-m}V/V
 \xrightarrow{F_V-1}J^{-pm}V/V
 \longrightarrow U_mR^1\Psi M\longrightarrow0.
 \tag{F5.1}
\]
In detail, a section with a pole of order $e>0$ has Frobenius pole
order $pe$, so its leading pole cannot cancel with its own order-$e$
term; regularity tests membership in these divisor lattices at
codimension-one points. This also proves the claimed kernel and
preimage assertions in (F5.1).

\paragraph{A pole cutoff independent of the coefficients.}
For $m\geq2$, the same quotient calculation as in Suzuki's
Proposition~3.1, using the projection formula for Frobenius and
$F_X^*V\simeq V$, gives
\[
 \operatorname{gr}_m R^1\Psi M
 \simeq B_Z\otimes_{\mathcal O_Z}V|_Z
                  \otimes_{\mathcal O_Z}J^{-m}|_Z,
 \qquad
 B_Z=(F_*\mathcal O_{pZ})/\mathcal O_Z.
 \tag{F5.2}
\]
The sheaf $B_Z$ is locally free on $Z$: the finite flat Frobenius
on the regular scheme $X$ makes $(F_*\mathcal O_X)/\mathcal O_X$
locally free, and the displayed sheaf is its base change to $Z$.

Fix a filtration of $\mathcal O_Z$ with successive terms line bundles
$P_\nu$ on components $D_{i(\nu)}$, obtained by removing one
component at a time from the effective cycle $Z$. Tensoring by the
locally free sheaves in (F5.2) gives successive terms
\[
 (B_Z|_{D_{i(\nu)}}\otimes P_\nu)
   \otimes L_a|_{D_{i(\nu)}}
   \otimes J^{-m}|_{D_{i(\nu)}}.
\]
For each of the finitely many fixed vector bundles
$B_Z|_{D_i}\otimes P_\nu$, choose an injection into a finite sum
of copies of a line bundle $H_\nu$ of positive degree. Such an
injection is obtained by globally generating a sufficiently ample
twist of the dual bundle and then dualizing. Since
$\deg L_a|_{D_i}=0$ and $\deg J|_{D_i}>0$, all the resulting line
bundles have negative degree for $m\geq m_0$, with one $m_0$
independent of $a,M,r$. Therefore
\[
 H^0(D,\operatorname{gr}_mR^1\Psi M)=0
 \quad(m\geq m_0).
\]
Filtered colimits and quasi-compactness now give
\[
 H^0(D,R^1\Psi M)=H^0(D,U_{m_0}R^1\Psi M),
 \tag{F5.3}
\]
increasing $m_0$ by one if necessary.

\paragraph{The remaining pole range has dimension at most a constant times rank.}
Suzuki's negative-degree filtration, Proposition~2.1(1), applies to
$J^{-pm}/\mathcal O_X$ and $J^{-pm}/J^{-m}$. Tensoring every layer
by a finite-order $L_a$ keeps its degree negative. Thus both
quotients tensored by $V$ have no global sections. Taking cohomology
of (F5.1) injects $H^0(D,U_mR^1\Psi M)$ into the kernel of
\[
 F_V-1:H^1(D,J^{-m}V/V)\longrightarrow
                    H^1(D,J^{-pm}V/V),
\]
and the map denoted $1$ here is injective. If a Frobenius-linear
map and an injective linear map have common domain of $k$-dimension
$N$, their difference has an $\mathbf F_p$-kernel of dimension at
most $N$: project to a linear splitting of the latter map and use
the fixed-space bound for a semilinear endomorphism. For completeness,
$\mathbf F_p$-linearly independent fixed vectors are $k$-linearly
independent: normalize a shortest $k$-linear relation to have first
coefficient $1$, apply Frobenius, and subtract; minimality forces
all coefficients into $\mathbf F_p$, a contradiction. This is the
quantitative form of Suzuki's Lemma~3.3.

A fixed filtration of $J^{-m_0}/\mathcal O_X$ has line-bundle layers
$T_\nu$ on smooth components $D_{i(\nu)}$. Their degrees are
negative. On a smooth projective curve $C$ of genus $g$, Riemann--Roch
therefore gives, for every degree-zero $L$,
\[
 h^1(C,T_\nu\otimes L)=g-1-\deg T_\nu.
\]
The finite filtration and the splitting of $V$ imply
\[
 \dim_k H^1(D,J^{-m_0}V/V)\leq C_1r
\]
for a constant $C_1$ fixed by $X,J,m_0$. Combining this with
(F5.3) gives
\[
 \dim_{\mathbf F_p}H^0(D,R^1\Psi M)\leq C_1r.
\]

Finally, the five-term nearby-cycle sequence and proper base change
bound $H^1(U,M)$ by $H^1(D,M)$ and $H^0(D,R^1\Psi M)$.
The Artin--Schreier sequence on the proper curve $D$ identifies
$H^1(D,M)$ with the Frobenius-fixed subgroup of $H^1(D,V|_D)$: on the finite-dimensional space of
global sections, $F_V-1$ is surjective over the algebraically closed
field. The normalization sequence of the nodal curve $D$ bounds
\[
 h^1(D,L_a|_D)\leq\sum_i g(D_i)+\#\{\text{nodes of }D\},
\]
since each restriction to a component has degree zero. Consequently
$\dim_{\mathbf F_p}H^1(D,M)\leq C_2r$. This yields the desired
bound with $C_A=C_1+C_2$ for the fixed cover.

For the last assertion, suppose $\pi_1(U)$ surjected onto the group
of F3 in characteristic $2$. Its modules $M_n$ have cyclic odd-order
monodromy. Inflation and F3b would give
\[
 n=\frac{\dim H^1(G,M_n)}{\dim M_n}
 \leq\frac{\dim H^1(U,M_n)}{\dim M_n}\leq C_A
\]
for every $n$, a contradiction.
\end{proof}

\paragraph{Exact scope of this advance.}
F5 bounds coefficients with abelian prime-to-$p$ monodromy and
excludes the specific metabelian obstruction F3 as a geometric
quotient. It does not bound all simple modules of a prosolvable
fundamental group: their finite monodromy can be nonabelian and
can have order divisible by $p$. It therefore does not yet prove
finite generation of the maximal prosolvable quotient, much less
the full local fundamental group. The fixed-cover argument depends
on the finite abelian inertia subgroup and does not extend formally
to those other monodromy groups.
